% 中期报告派生正文；原学位论文章节、实验数据与图像保持只读。
% 核心主张逐项核对：docs/middle-defense/magnetic-evidence.tsv。
\chapter{磁感知与主动横切反演的阶段结果}
\label{chap:middle-magnetic}

海缆磁感知需要同时回答三个问题：如何在存在先验偏差时找到磁异常，如何从异常剖面获得可解释的几何量，以及几何量不足时是否仍能开展有限巡检。本阶段在原有磁场建模与电缆跟踪接口基础上，完成独立的第二版处理链，将先验引导搜索、磁检测、主动横切、半高全宽估计和沿缆跟踪连接起来。研究重点由单次反演误差扩展为观测形成、有效性检查与任务许可之间的关系。本章先给出几何定义，再用受控剖面比较和原生动力学闭环分别检验估计稳定性与任务可达性。

\section{磁观测量与反演几何定义}
\label{sec:middle-mag-geometry}

通电导线的磁场由电流分布和观测几何共同决定，其基本关系来自毕奥--萨伐尔定律\cite{griffiths2017electrodynamics}。对于三相海缆，芯线相位、绞合结构与铠装会影响可见场强和空间分布。采用绝对幅值反推距离时，未知电流与整体衰减会直接进入反演；因此，本阶段考察能否利用横切剖面的几何宽度降低对绝对幅值的依赖。该思路以声明的场形为前提，不将单直导线和多芯海缆的空间衰减混为一谈。

三轴采样先经标定与同频提取，记目标频率上的三轴均方根幅值为 $R_1,R_2,R_3$，合成观测量为
\begin{equation}
 B=\sqrt{R_1^2+R_2^2+R_3^2}.
 \label{eq:middle-mag-envelope}
\end{equation}
在三轴增益与正交性已经校正、测量点不变的条件下，该模量不因坐标旋转而改变；它并不消除杆臂运动、通道失配及沿缆场变化。后文受控仿真直接使用标量包络，以便单独分析剖面几何；三轴采集链的实际性能在平台验证部分讨论。

建立局部坐标：$s$ 沿电缆切向，$x$ 为横切方向，$c$ 为电缆在横切坐标中的中心，$z>0$ 为磁传感器与电缆之间的垂直距离。若海床在局部可近似为平面，传感器离底高度为 $h$，则埋深为 $d=z-h$。这里的 $z$ 既不是艇体水深，也不是非零横偏处的斜距；由 $z$ 得到 $d$ 还需要同步高度和统一的传感器空间基准。船体导航位置应经杆臂改正后用于横切坐标计算。

在对称、单峰的解析包络模型下，含局部背景的响应可表示为
\begin{equation}
 B(x)=b+\frac{C}{[(x-c)^2+z^2]^{p/2}},\qquad C>0,
 \label{eq:middle-mag-profile}
\end{equation}
其中 $b$ 是局部标量背景，$C$ 汇集幅值尺度，$p$ 表征径向衰减。设峰值为 $B_{\mathrm{pk}}$，半高水平应取 $b+(B_{\mathrm{pk}}-b)/2$，而非含背景峰值的二分之一。左右交点记为 $x_L,x_R$，空间半高全宽（full width at half maximum，FWHM）及交点中点为
\begin{equation}
 W_x=x_R-x_L,\qquad
 \hat c=\frac{x_L+x_R}{2},\qquad
 z=\kappa_p W_x,\quad
 \kappa_p=\frac{1}{2\sqrt{2^{2/p}-1}}.
 \label{eq:middle-mag-width}
\end{equation}
本文二维线偶极近似采用 $p=2$ 的标量包络，因此 $z=W_x/2$；它不同于三维点偶极模型。无限长单直导线的模量包络对应 $p=1$，因此 $z=W_x/(2\sqrt{3})$。整体增益只改变峰高，在模型与背景条件成立时不改变半高交点，但错误的 $p$ 会系统性改变距离换算。有限导线及实测海缆应另行标定宽度与距离的关系。

横向中心和垂向尺度具有不同的可辨识条件。更一般地，对 $B(x)=b+A g((x-c)/z)$，若 $g$ 为偶函数且在正半轴严格递减，交点中点仍为 $c$；宽度则为 $2zq_g$，其中 $g(q_g)=g(0)/2$。因而，未知对称场形下的中心可以有意义，宽度到距离的系数却未必已知。这一分解成为后续分别管理横向定位与埋深输出的依据，同时也指出融合邻源破坏单源解释时需要额外的目标关联信息。

\section{从先验搜索到主动横切的处理流程}
\label{sec:middle-mag-workflow}

原有近缆反演流程默认能够持续观测到目标，难以解释初始位置偏差较大时怎样进入该条件。第二版将搜索作为独立阶段：以电缆先验折线、初始位置不确定度和授权边界生成规则往返测线。搜索宽度至少覆盖所声明的位置不确定范围，测线间距决定空间覆盖密度，航速决定单位距离上的采样密度及搜索时间。规划器采用可解释的平行或横贯测线；若所需走廊超出边界则拒绝该计划。搜索是否找到异常，与后续能否反演几何分别评价，从而避免将已在缆附近的局部演示误写成完整搜索。

磁检测器维护滑动窗口，以历史低分位数估计背景，以差分序列的中位数绝对偏差估计噪声，并设置正的噪声下限。检测分数是当前正异常相对噪声的比值；进入检测和退出检测使用不同阈值，并要求连续样本满足条件。迟滞和持续计数抑制临界值附近的反复切换，峰值趋势作为补充证据。这里的分数和后续质量分是工程诊断量，未被标定为真实检测概率。单次高分只触发进一步观察，不直接授权埋深或目标跟踪。

确认阶段通过横切机动主动产生空间变化，并用导航位置构造 $B(x)$。均速直线且跨切角 $\alpha$ 固定时，可以用 $W_x=v|\sin\alpha|\Delta t$ 将时间宽度换算为空间宽度；存在加减速、转弯或横流时，本实现优先使用实测空间坐标与磁样本配对。对单调横切腿完成排序、重复位置合并和局部平滑后，显式寻找峰值两侧的半高交点并插值，输出峰值、背景、交点、宽度、中点以及诊断原因。观测不足时保留已测剖面，不通过外推补造缺失交点。

扫描幅度与扫描方式均服务于双侧覆盖。过窄扫描可能只看到峰值附近，无法区分背景与场形；无条件增幅又会增加沿缆漂移、转向负担和接入距离。因此采用受边界与运动能力限制的逐级扩幅，并检查峰值包围、双侧采样、对称性、信噪比与模型残差。曲率和横流条件变化时，局部浅角横切或连续之字形可以提供多站观测，但有效磁剖面仍须与可执行接入轨迹同时成立。方法选择不能只以 FWHM 有效率最大为目标。

处理链将磁异常、横向中心、埋深和执行许可分别保存。完整横切支持横向几何后，系统依据估计线形成航向和横向修正，通过既有控制器接入有限沿缆路线；缺少场形标定或同步离底高度时，埋深保持空值。横向观测设置有效范围与时效，导航、目标关联或观测质量失效时进入暂停状态并请求后续处置。这使任务能够利用已有可信信息，也使不满足条件的量在接口层可被明确识别。

\section{稳定性比较与场形失配的拒绝}
\label{sec:middle-mag-stability}

稳定性实验比较显式空间 FWHM、自由幅值剖面拟合、标定幅值反演和未扣背景半高四种方法。七类场景包括标称、增益、背景、噪声、扫描偏心、组合扰动和场形失配；每场景每方法使用 200 个固定种子，共 5600 次方法运行。距离真值从 $2.5$--$5.0~\mathrm{m}$ 取样，每条剖面为 241 点。这一冻结基线产生于新增场形判别门之前，后续修复没有覆盖它。

比较的输入条件需要明确：FWHM 接收生成器给出的标量背景和噪声标准差，标定幅值法也扣除该标量背景；叠加的空间背景斜率仍进入观测。自由拟合同时估计背景、峰高、中心和尺度。幅值法采用名义源强，因此增益变化会形成标定失配。本组用于比较受控条件下的几何反演和发布行为，不单独评价在线背景估计性能。误差统计仅针对有效发布结果；有效率及大误差发布比例以该组全部 200 次为分母，拒绝样本不以零误差计入。

图~\ref{fig:middle-mag-stability}表明，方法收益随扰动类型变化。标称条件下，FWHM、自由拟合和幅值法的平均绝对误差（mean absolute error，MAE）分别为 $0.017$、$0.007$ 和 $0.004~\mathrm{m}$，FWHM 并非标称精度最优。增益变化时，FWHM 与幅值法的 MAE 分别为 $0.027$ 和 $0.772~\mathrm{m}$，体现了宽度消去整体幅值的作用。组合扰动下，FWHM 的发布率为 $97.0\%$，发布 MAE 为 $0.105~\mathrm{m}$；自由拟合和幅值法均全部发布，MAE 分别为 $0.171$ 和 $0.678~\mathrm{m}$。三者误差超过 $0.5~\mathrm{m}$ 仍发布的比例分别为 $5.5\%$、$8.5\%$ 和 $61.0\%$。

\begin{figure}[htbp]
 \centering
 \includegraphics[width=15cm]{middle/figures/magnetic_stability.pdf}
 \caption{受控剖面的四方法稳定性比较。每场景每方法 200 次，误差针对发布结果，发布率与大误差发布比例按全部样本计算；箱线图省略离群点，大误差仍计入风险比例。}
 \label{fig:middle-mag-stability}
\end{figure}

扫描偏心场景中，FWHM 只发布 $74.0\%$ 的剖面，且没有大于 $0.5~\mathrm{m}$ 的误差发布，说明质量检查以覆盖不足时的拒绝换取输出完整性。但该机制不能发现所有物理失配：真实剖面改为单直导线而继续采用偶极系数时，FWHM 全部有效，MAE 却达到 $2.753~\mathrm{m}$，200 次误差均超过 $0.5~\mathrm{m}$；自由拟合的 MAE 也达到 $1.483~\mathrm{m}$。峰值、交点和较小拟合残差并不自动证明距离正确，这一反例直接促成了后续的场形检查。

场形判别保持实测峰值、背景、中点和宽度不变，在相同自由度下比较声明模型与备选模型的归一化均方根残差。只有备选残差比声明残差至少低 $0.02$，且小于后者的 $0.75$ 倍时，才判定备选场形明显更优。另要求从交点中点到扫描两端的较短距离不小于 $0.65W_x$，以保证峰外形状信息充分。任一条件触发拒绝时，仍保留宽度和诊断，不发布由错误模型换算的距离。

阈值探索使用种子 0--199，图~\ref{fig:middle-mag-shape}采用独立种子 200--399 的留出结果。真实单直导线误声明为偶极、真实偶极误声明为单直导线两组中，旧门分别放行全部 200 次大误差，新门将这 400 次均拒绝。正确声明的直导线、标称偶极、增益变化与背景变化组维持全部有效；噪声组由 $100\%$ 降为 $99\%$，扫描偏心由 $74.0\%$ 降为 $64.5\%$，组合扰动由 $92.5\%$ 降为 $87.5\%$。组合扰动的大误差发布比例仍为 $3.0\%$。因此，新增门针对已测试场形混淆有效，同时付出误拒与覆盖要求提高的代价。

\begin{figure}[htbp]
 \centering
 \includegraphics[width=15cm]{middle/figures/magnetic_shape_guard.pdf}
 \caption{独立留出集的场形失配风险与发布代价。上图为两类场形相互误声明的 400 条剖面，下图合并展示正确声明模型在常规和困难条件下的发布率；大误差指距离误差超过 $0.5~\mathrm{m}$。}
 \label{fig:middle-mag-shape}
\end{figure}

形状比较仍不能识别保持归一化形状的系数漂移，也不能穷尽未知多芯、多源场。由此形成的方案调整是：距离发布需要独立几何真值支持的受控横切标定，训练扫描与留出扫描分离，并检查适用距离范围；横向巡检则按其自身观测条件单独评估。模型适用性、真实同步高度和目标身份的实物验证进入后续工作，不用更高的拟合自由度替代必要测量。

\section{有限路线巡检与条件性监督结果}
\label{sec:middle-mag-closure}

任务验证首先采用磁算法、既有 PID、72 字节协议编解码和原生 PVS 动力学组成的独立回路。目标为已知孤立直线磁源，理想导航、零横流，先验线与真实线横向相差 $6~\mathrm{m}$。单直导线与偶极标量包络分别取峰值 $20~\mathrm{nT}$、背景 $1~\mathrm{nT}$、垂向尺度 $3~\mathrm{m}$，加标准差 $0.03~\mathrm{nT}$ 的高斯噪声。名义速度 $0.3~\mathrm{m/s}$，决策步长 $0.25~\mathrm{s}$，动力学积分步长 $0.05~\mathrm{s}$。本组固定二维磁包络，以分离主动观测、接入与航迹控制的作用。

浅角横切从 $(-250,-15)~\mathrm{m}$ 指向 $(-80,15)~\mathrm{m}$，之后以 $20~\mathrm{m}$ 前向距离接入估计线。名义任务为 $80~\mathrm{m}$，两端各留 $5~\mathrm{m}$ 缓冲，核验核心为沿缆坐标 $[5,75]~\mathrm{m}$。两种场形各三个种子完成核心覆盖，三个无源对照均拒绝进入跟踪；全部九次埋深为空。核心长度在首个种子观察后、另外两个种子复验前确定，因此本组报告为受限条件下的复验，不作为全独立海况抽样。

\begin{figure}[htbp]
 \centering
 \includegraphics[width=15cm]{middle/figures/magnetic_mission_process.pdf}
 \caption{限定条件下的主动磁巡检。地图和中心估计采用种子 101，横偏包含两场形各三个种子；阴影为 $70~\mathrm{m}$ 核心段，全部核心跟踪样本含回归修正均计入。地图横纵轴比例不同，埋深始终未发布。}
 \label{fig:middle-mag-process}
\end{figure}

图~\ref{fig:middle-mag-process}将位置估计与实际运动分开展示：磁中心接近真实线后，载体仍需经历接入和横向调节。以全部核心跟踪样本计算，六条正例的真实横偏 MAE 为 $0.771$--$0.774~\mathrm{m}$，最大值为 $1.321~\mathrm{m}$，满足该试验预设的均值不超过 $1~\mathrm{m}$、最大值不超过 $2~\mathrm{m}$。评价包含回归修正阶段，未用业务阶段标签剔除较大误差。该结果支持有限横向任务的可达性，也说明不能用磁中心的小误差代替载体跟踪误差。

进一步的留出实验加入扰流、导航误差与相干邻源。在原处理链之外，独立监督层用报告位移与世界系地速积分的一致性检查跳变和漂移，用冻结的完整横切剖面检查跟踪阶段的幅值变化。接入仅作一项有条件调整：测得的反向扰动代理超过 $0.02~\mathrm{m/s}$ 时，将前向接入距离由 $20~\mathrm{m}$ 改为 $5~\mathrm{m}$；该代理包含侧滑，不当作直接海流测量。

指定目标授权额外使用独立外部参考区间 $y=6\pm0.05~\mathrm{m}$，并假设实际水平导航误差不超过 $0.10~\mathrm{m}$，磁中心误差预算另设为 $0.10~\mathrm{m}$。若参考中心为 $c_{\mathrm{ref}}$、半宽为 $\epsilon_{\mathrm{ref}}$，报告位置为 $\hat y$，则条件成立时有
\begin{equation}
 |y-c|\leq |\hat y-c_{\mathrm{ref}}|
       +\epsilon_{\mathrm{nav}}+\epsilon_{\mathrm{ref}}.
 \label{eq:middle-mag-bound}
\end{equation}
该上界将真实误差约束连接到可检查的报告量；参考与导航误差界必须实际成立，不能由算法自报置信度证明。外部参考给监督组提供了额外信息，因此两方法比较同时考察监督与接入的作用，不解释为等信息条件下的磁反演精度排名。

冻结参数后，20 种条件、三个留出种子与两种方法形成 120 次闭环。九类声明正例包括孤立源、未知对称包络、正反小稳态流、周期流、小导航扰动、弱邻源与组合扰动，监督组完成 $27/27$ 次；完整任务的全核心真实横偏 MAE 为 $0.426$--$0.816~\mathrm{m}$，最大横偏范围为 $0.775$--$1.552~\mathrm{m}$。其余十一类对照撤权 $33/33$ 次，监督组错误几何授权与错误任务完成均为 $0/60$。拒绝不计为巡检完成，拒绝前的截短误差也不与完整任务排名。

\begin{figure}[htbp]
 \centering
 \includegraphics[width=15cm]{middle/figures/magnetic_supervision.pdf}
 \caption{接入修正与目标授权的留出对照（种子 501）。监督组含外部参考；左图灰区为旧链未覆盖路段，右图为旧链误差按对应阈值归一化，监督组在该融合邻源案例中拒绝目标几何授权。}
 \label{fig:middle-mag-supervision}
\end{figure}

旧链到端 $39/60$ 次，但存在 $12/60$ 次错误几何授权。图~\ref{fig:middle-mag-supervision}中的融合邻源案例表明，载体路线误差合格与磁中心属于指定目标是两项不同结论。监督组拒绝不一致的目标关联；无外部参考的九次消融虽均结束几何路线，指定目标完成仍为 $0/9$。仅修改接入、不增加监督的三个反向小流案例完成任务，进一步区分了运动几何收益和目标参考收益。

反向阶跃、超界流等条件仍未完成巡检。监督撤权后持续记录 $60~\mathrm{s}$ 零命令，并在 $30~\mathrm{s}$ 请求上级恢复，原生动力学仍产生滑行或漂移；本组没有模拟上级恢复动作。由此，本阶段已形成主动磁观测、显式质量检查和条件性横向任务的仿真证据。未知场形埋深、真实目标关联、导航误差界与水中恢复能力分别需要后续标定和实物验证，历史 ROS 链路结果与本次独立回路保持各自的版本和验证层级。
